% Options for packages loaded elsewhere
\PassOptionsToPackage{unicode}{hyperref}
\PassOptionsToPackage{hyphens}{url}
\PassOptionsToPackage{dvipsnames,svgnames,x11names}{xcolor}
%
\documentclass[
  11pt,
]{article}

\usepackage{amsmath,amssymb}
\usepackage{iftex}
\ifPDFTeX
  \usepackage[T1]{fontenc}
  \usepackage[utf8]{inputenc}
  \usepackage{textcomp} % provide euro and other symbols
\else % if luatex or xetex
  \usepackage{unicode-math}
  \defaultfontfeatures{Scale=MatchLowercase}
  \defaultfontfeatures[\rmfamily]{Ligatures=TeX,Scale=1}
\fi
\usepackage{lmodern}
\ifPDFTeX\else  
    % xetex/luatex font selection
\fi
% Use upquote if available, for straight quotes in verbatim environments
\IfFileExists{upquote.sty}{\usepackage{upquote}}{}
\IfFileExists{microtype.sty}{% use microtype if available
  \usepackage[]{microtype}
  \UseMicrotypeSet[protrusion]{basicmath} % disable protrusion for tt fonts
}{}
\makeatletter
\@ifundefined{KOMAClassName}{% if non-KOMA class
  \IfFileExists{parskip.sty}{%
    \usepackage{parskip}
  }{% else
    \setlength{\parindent}{0pt}
    \setlength{\parskip}{6pt plus 2pt minus 1pt}}
}{% if KOMA class
  \KOMAoptions{parskip=half}}
\makeatother
\usepackage{xcolor}
\usepackage[margin=1in]{geometry}
\setlength{\emergencystretch}{3em} % prevent overfull lines
\setcounter{secnumdepth}{5}
% Make \paragraph and \subparagraph free-standing
\ifx\paragraph\undefined\else
  \let\oldparagraph\paragraph
  \renewcommand{\paragraph}[1]{\oldparagraph{#1}\mbox{}}
\fi
\ifx\subparagraph\undefined\else
  \let\oldsubparagraph\subparagraph
  \renewcommand{\subparagraph}[1]{\oldsubparagraph{#1}\mbox{}}
\fi

\usepackage{color}
\usepackage{fancyvrb}
\newcommand{\VerbBar}{|}
\newcommand{\VERB}{\Verb[commandchars=\\\{\}]}
\DefineVerbatimEnvironment{Highlighting}{Verbatim}{commandchars=\\\{\}}
% Add ',fontsize=\small' for more characters per line
\usepackage{framed}
\definecolor{shadecolor}{RGB}{241,243,245}
\newenvironment{Shaded}{\begin{snugshade}}{\end{snugshade}}
\newcommand{\AlertTok}[1]{\textcolor[rgb]{0.68,0.00,0.00}{#1}}
\newcommand{\AnnotationTok}[1]{\textcolor[rgb]{0.37,0.37,0.37}{#1}}
\newcommand{\AttributeTok}[1]{\textcolor[rgb]{0.40,0.45,0.13}{#1}}
\newcommand{\BaseNTok}[1]{\textcolor[rgb]{0.68,0.00,0.00}{#1}}
\newcommand{\BuiltInTok}[1]{\textcolor[rgb]{0.00,0.23,0.31}{#1}}
\newcommand{\CharTok}[1]{\textcolor[rgb]{0.13,0.47,0.30}{#1}}
\newcommand{\CommentTok}[1]{\textcolor[rgb]{0.37,0.37,0.37}{#1}}
\newcommand{\CommentVarTok}[1]{\textcolor[rgb]{0.37,0.37,0.37}{\textit{#1}}}
\newcommand{\ConstantTok}[1]{\textcolor[rgb]{0.56,0.35,0.01}{#1}}
\newcommand{\ControlFlowTok}[1]{\textcolor[rgb]{0.00,0.23,0.31}{#1}}
\newcommand{\DataTypeTok}[1]{\textcolor[rgb]{0.68,0.00,0.00}{#1}}
\newcommand{\DecValTok}[1]{\textcolor[rgb]{0.68,0.00,0.00}{#1}}
\newcommand{\DocumentationTok}[1]{\textcolor[rgb]{0.37,0.37,0.37}{\textit{#1}}}
\newcommand{\ErrorTok}[1]{\textcolor[rgb]{0.68,0.00,0.00}{#1}}
\newcommand{\ExtensionTok}[1]{\textcolor[rgb]{0.00,0.23,0.31}{#1}}
\newcommand{\FloatTok}[1]{\textcolor[rgb]{0.68,0.00,0.00}{#1}}
\newcommand{\FunctionTok}[1]{\textcolor[rgb]{0.28,0.35,0.67}{#1}}
\newcommand{\ImportTok}[1]{\textcolor[rgb]{0.00,0.46,0.62}{#1}}
\newcommand{\InformationTok}[1]{\textcolor[rgb]{0.37,0.37,0.37}{#1}}
\newcommand{\KeywordTok}[1]{\textcolor[rgb]{0.00,0.23,0.31}{#1}}
\newcommand{\NormalTok}[1]{\textcolor[rgb]{0.00,0.23,0.31}{#1}}
\newcommand{\OperatorTok}[1]{\textcolor[rgb]{0.37,0.37,0.37}{#1}}
\newcommand{\OtherTok}[1]{\textcolor[rgb]{0.00,0.23,0.31}{#1}}
\newcommand{\PreprocessorTok}[1]{\textcolor[rgb]{0.68,0.00,0.00}{#1}}
\newcommand{\RegionMarkerTok}[1]{\textcolor[rgb]{0.00,0.23,0.31}{#1}}
\newcommand{\SpecialCharTok}[1]{\textcolor[rgb]{0.37,0.37,0.37}{#1}}
\newcommand{\SpecialStringTok}[1]{\textcolor[rgb]{0.13,0.47,0.30}{#1}}
\newcommand{\StringTok}[1]{\textcolor[rgb]{0.13,0.47,0.30}{#1}}
\newcommand{\VariableTok}[1]{\textcolor[rgb]{0.07,0.07,0.07}{#1}}
\newcommand{\VerbatimStringTok}[1]{\textcolor[rgb]{0.13,0.47,0.30}{#1}}
\newcommand{\WarningTok}[1]{\textcolor[rgb]{0.37,0.37,0.37}{\textit{#1}}}

\providecommand{\tightlist}{%
  \setlength{\itemsep}{0pt}\setlength{\parskip}{0pt}}\usepackage{longtable,booktabs,array}
\usepackage{calc} % for calculating minipage widths
% Correct order of tables after \paragraph or \subparagraph
\usepackage{etoolbox}
\makeatletter
\patchcmd\longtable{\par}{\if@noskipsec\mbox{}\fi\par}{}{}
\makeatother
% Allow footnotes in longtable head/foot
\IfFileExists{footnotehyper.sty}{\usepackage{footnotehyper}}{\usepackage{footnote}}
\makesavenoteenv{longtable}
\usepackage{graphicx}
\makeatletter
\def\maxwidth{\ifdim\Gin@nat@width>\linewidth\linewidth\else\Gin@nat@width\fi}
\def\maxheight{\ifdim\Gin@nat@height>\textheight\textheight\else\Gin@nat@height\fi}
\makeatother
% Scale images if necessary, so that they will not overflow the page
% margins by default, and it is still possible to overwrite the defaults
% using explicit options in \includegraphics[width, height, ...]{}
\setkeys{Gin}{width=\maxwidth,height=\maxheight,keepaspectratio}
% Set default figure placement to htbp
\makeatletter
\def\fps@figure{htbp}
\makeatother

\usepackage{algorithm}
\usepackage{algorithmic}
\usepackage{amsmath}
\usepackage{float}
\usepackage{booktabs}
\floatname{algorithm}{Algorithm}
\makeatletter
\makeatother
\makeatletter
\makeatother
\makeatletter
\@ifpackageloaded{caption}{}{\usepackage{caption}}
\AtBeginDocument{%
\ifdefined\contentsname
  \renewcommand*\contentsname{Table of contents}
\else
  \newcommand\contentsname{Table of contents}
\fi
\ifdefined\listfigurename
  \renewcommand*\listfigurename{List of Figures}
\else
  \newcommand\listfigurename{List of Figures}
\fi
\ifdefined\listtablename
  \renewcommand*\listtablename{List of Tables}
\else
  \newcommand\listtablename{List of Tables}
\fi
\ifdefined\figurename
  \renewcommand*\figurename{Figure}
\else
  \newcommand\figurename{Figure}
\fi
\ifdefined\tablename
  \renewcommand*\tablename{Table}
\else
  \newcommand\tablename{Table}
\fi
}
\@ifpackageloaded{float}{}{\usepackage{float}}
\floatstyle{ruled}
\@ifundefined{c@chapter}{\newfloat{codelisting}{h}{lop}}{\newfloat{codelisting}{h}{lop}[chapter]}
\floatname{codelisting}{Listing}
\newcommand*\listoflistings{\listof{codelisting}{List of Listings}}
\makeatother
\makeatletter
\@ifpackageloaded{caption}{}{\usepackage{caption}}
\@ifpackageloaded{subcaption}{}{\usepackage{subcaption}}
\makeatother
\makeatletter
\@ifpackageloaded{tcolorbox}{}{\usepackage[skins,breakable]{tcolorbox}}
\makeatother
\makeatletter
\@ifundefined{shadecolor}{\definecolor{shadecolor}{rgb}{.97, .97, .97}}
\makeatother
\makeatletter
\makeatother
\makeatletter
\makeatother
\ifLuaTeX
  \usepackage{selnolig}  % disable illegal ligatures
\fi
\IfFileExists{bookmark.sty}{\usepackage{bookmark}}{\usepackage{hyperref}}
\IfFileExists{xurl.sty}{\usepackage{xurl}}{} % add URL line breaks if available
\urlstyle{same} % disable monospaced font for URLs
\hypersetup{
  pdftitle={Stimulus-Locked Turn Rate Analysis Pipeline - Pseudocode Guide},
  pdfauthor={INDYsim Analysis Pipeline},
  colorlinks=true,
  linkcolor={blue},
  filecolor={Maroon},
  citecolor={Blue},
  urlcolor={Blue},
  pdfcreator={LaTeX via pandoc}}

\title{Stimulus-Locked Turn Rate Analysis Pipeline - Pseudocode Guide}
\author{INDYsim Analysis Pipeline}
\date{Invalid Date}

\begin{document}
\maketitle
\ifdefined\Shaded\renewenvironment{Shaded}{\begin{tcolorbox}[sharp corners, breakable, enhanced, boxrule=0pt, interior hidden, frame hidden, borderline west={3pt}{0pt}{shadecolor}]}{\end{tcolorbox}}\fi

\renewcommand*\contentsname{Table of contents}
{
\hypersetup{linkcolor=}
\setcounter{tocdepth}{3}
\tableofcontents
}
\newpage
\tableofcontents
\newpage

\hypertarget{overview}{%
\section{Overview}\label{overview}}

This document provides a detailed pseudocode guide for the
stimulus-locked turn rate analysis pipeline. The pipeline processes H5
trajectory files, performs MAGAT segmentation, detects behavioral
events, and calculates stimulus-locked turn rates.

\hypertarget{variable-reference}{%
\section{Variable Reference}\label{variable-reference}}

This section provides a comprehensive reference table of all variables
used throughout the pipeline.

\begin{table}[H]
\centering
\footnotesize
\begin{tabular}{p{3cm}p{2cm}p{9cm}}
\hline
\textbf{Variable} & \textbf{Type} & \textbf{Description} \\
\hline
\multicolumn{3}{l}{\textbf{File and Data Structures}} \\
\hline
$h5\_file$ & str & Path to H5 trajectory file \\
$h5\_data$ & dict & Loaded H5 file data structure \\
$fps$ & float & Frames per second (frame rate) \\
$stimulus\_df$ & DataFrame & Stimulus timing data \\
$track\_data$ & dict & Raw track data from H5 file \\
$df$ & DataFrame & Trajectory DataFrame with features \\
$traj\_df$ & DataFrame & Trajectory DataFrame (alias for $df$) \\
$aligned\_df$ & DataFrame & Trajectory aligned with stimulus timing \\
$events\_df$ & DataFrame & Event records DataFrame (from CSV) \\
$trajectories\_df$ & DataFrame & Trajectories DataFrame (from CSV) \\
$event\_records$ & DataFrame & Binned event records for single track \\
\hline
\multicolumn{3}{l}{\textbf{Trajectory Features}} \\
\hline
$x, y$ & array & Position coordinates (cm or mm) \\
$heading$ & array & Heading angle (radians, $[0, 2\pi)$) \\
$speed$ & array & Instantaneous speed (mm/s or cm/s) \\
$curvature$ & array & Trajectory curvature (rad/cm) \\
$\Delta\theta$ & array & Heading change (radians, wrapped to $[-\pi, +\pi]$) \\
$accel$ & array & Acceleration (mm/s$^2$ or cm/s$^2$) \\
$dt$ & array & Time differences between frames (seconds) \\
$time$ & array & Time array (seconds) \\
$n\_frames$ & int & Total number of frames in track \\
\hline
\multicolumn{3}{l}{\textbf{Behavioral States (Boolean Arrays)}} \\
\hline
$is\_run$ & array[bool] & True during run periods \\
$is\_reorientation$ & array[bool] & True during reorientation periods \\
$is\_turn$ & array[bool] & True during turn periods (reorientations with pauses) \\
$is\_pause$ & array[bool] & True during pause periods \\
$is\_head\_swing$ & array[bool] & True during head swing periods \\
$is\_reversal$ & array[bool] & True during reversal periods \\
\hline
\multicolumn{3}{l}{\textbf{START Event Detection}} \\
\hline
$is\_reorientation\_start$ & array[bool] & True at reorientation START events (False$\rightarrow$True transitions) \\
$is\_turn\_start$ & array[bool] & True at turn START events \\
$is\_pause\_start$ & array[bool] & True at pause START events \\
$is\_reo\_padded$ & array[bool] & Padded boolean array for transition detection \\
$reo\_start\_mask$ & array[bool] & Mask identifying START transitions \\
\hline
\multicolumn{3}{l}{\textbf{MAGAT Segmentation}} \\
\hline
$runs$ & list & List of run intervals: $[(start\_idx, end\_idx), ...]$ \\
$head\_swings$ & list & List of head swing intervals \\
$reorientations$ & list & List of reorientation intervals \\
$turns$ & list & List of turn intervals (reorientations with pauses) \\
$segmentation$ & dict & Complete segmentation dictionary \\
$magat\_segmentation$ & dict & MAGAT segmentation results \\
\hline
\multicolumn{3}{l}{\textbf{MAGAT Thresholds}} \\
\hline
$curv\_cut$ & float & Curvature threshold (default: 0.4 rad/cm) \\
$theta\_cut$ & float & Body bend threshold (default: $\pi/2$ rad) \\
$stop\_speed\_cut$ & float & Speed threshold for run end (default: 2.0 mm/s) \\
$start\_speed\_cut$ & float & Speed threshold for run start (default: 3.0 mm/s) \\
$minRunTime$ & float & Minimum run duration (default: 2.5s) \\
$minRunLength$ & float & Minimum run path length (optional) \\
$minRunSpeed$ & float & Minimum average run speed (optional) \\
\hline
\multicolumn{3}{l}{\textbf{MAGAT Intermediate Variables}} \\
\hline
$body\_theta$ & array & Body bend angle (radians) \\
$vel\_dp$ & array & Velocity dot product (head alignment) \\
$high\_curv$ & array[bool] & High curvature condition \\
$head\_swinging$ & array[bool] & Head swinging condition \\
$speed\_low$ & array[bool] & Low speed condition \\
$not\_a\_run$ & array[bool] & Not-a-run condition (OR of end conditions) \\
$speed\_high$ & array[bool] & High speed condition \\
$head\_aligned$ & array[bool] & Head aligned condition \\
$is\_a\_run$ & array[bool] & Is-a-run condition (AND of start conditions) \\
$head\_swing\_start$ & array[bool] & Head swing start condition \\
$head\_swing\_stop$ & array[bool] & Head swing stop condition \\
\hline
\multicolumn{3}{l}{\textbf{Pause Detection}} \\
\hline
$speed\_threshold$ & float & Speed threshold for pause (default: 0.001 mm/s) \\
$min\_duration$ & float & Minimum pause duration (default: 0.2s) \\
$pause\_starts$ & list & List of pause start frame indices \\
$pause\_ends$ & list & List of pause end frame indices \\
$duration$ & float & Duration of a pause (seconds) \\
\hline
\multicolumn{3}{l}{\textbf{Turn Extraction}} \\
\hline
$reo\_start, reo\_end$ & int & Reorientation start/end frame indices \\
$reo\_pause\_frames$ & array[bool] & Pause frames within a reorientation \\
\hline
\multicolumn{3}{l}{\textbf{Stimulus Alignment}} \\
\hline
$led1Val$ & array & LED1 value (interpolated) \\
$led12Val\_ton$ & array & Time within LED ON period \\
$led12Val\_toff$ & array & Time within LED OFF period \\
$stimulus\_on$ & array[bool] & True when stimulus is ON \\
$stimulus\_onset$ & array[bool] & True at stimulus onset transitions \\
$stim\_idx$ & int & Index of nearest stimulus time \\
$threshold$ & float & LED value threshold for stimulus detection \\
\hline
\multicolumn{3}{l}{\textbf{Event Record Creation}} \\
\hline
$track\_id$ & int & Track identifier \\
$exp\_id$ & int & Experiment identifier \\
$experiment\_id$ & str & Experiment ID string \\
$bin\_size$ & float & Bin size for event records (default: 0.05s) \\
$b$ & int & Bin index \\
$bin\_start$ & float & Start time of bin (seconds) \\
$bin\_end$ & float & End time of bin (seconds) \\
$bin\_data$ & DataFrame & Data within bin \\
\hline
\multicolumn{3}{l}{\textbf{Cycle Extraction}} \\
\hline
$cycles$ & list & List of cycle dictionaries \\
$c$ & dict & Single cycle dictionary \\
$cycle\_num$ & int & Cycle number \\
$onset\_frame$ & int & Stimulus onset frame index \\
$onset\_time$ & float & Stimulus onset time (t=0 reference, seconds) \\
$led\_on\_frame$ & int & LED ON frame index \\
$led\_on\_time$ & float & LED ON time (seconds) \\
$led\_off\_frame$ & int & LED OFF frame index \\
$led\_off\_time$ & float & LED OFF time (seconds) \\
$pulse\_duration$ & float & Pulse duration (seconds) \\
$eti$ & float & Elapsed time index (OFF to next onset, seconds) \\
$baseline\_start\_time$ & float & Baseline start time (10s before onset) \\
$cycle\_start\_time$ & float & Cycle start time (analysis window start) \\
$cycle\_end\_time$ & float & Cycle end time (analysis window end) \\
\hline
\multicolumn{3}{l}{\textbf{Per-Cycle Statistics}} \\
\hline
$cycle\_start$ & float & Cycle start time (seconds) \\
$cycle\_end$ & float & Cycle end time (seconds) \\
$reo\_starts$ & DataFrame & Reorientation start events \\
$turn\_starts$ & DataFrame & Turn start events \\
$cycle\_reos$ & DataFrame & Reorientations within cycle \\
$n\_cycle\_reos$ & int & Number of reorientations in cycle \\
$cycle\_turns$ & DataFrame & Turns within cycle \\
$n\_cycle\_turns$ & int & Number of turns in cycle \\
$cycle\_duration\_min$ & float & Cycle duration in minutes \\
$turn\_rate$ & float & Turn rate for cycle (min$^{-1}$) \\
$n\_bins$ & int & Number of bins in cycle \\
$bin\_reos$ & DataFrame & Reorientations within bin \\
$n\_bin\_reos$ & int & Number of reorientations in bin \\
$n\_bin\_turns$ & int & Number of turns in bin \\
$bin\_duration\_min$ & float & Bin duration in minutes \\
$bin\_rate$ & float & Turn rate for bin (min$^{-1}$) \\
$time\_rel\_onset$ & float & Time relative to stimulus onset (seconds) \\
\hline
\multicolumn{3}{l}{\textbf{Stimulus-Locked Turn Rate}} \\
\hline
$bin\_edges$ & array & Bin edge times (seconds, relative to onset) \\
$t\_reo$ & float & Reorientation start time (seconds) \\
$t\_rel$ & float & Time relative to onset (seconds) \\
$bin\_idx$ & int & Bin index \\
$counts$ & array & Event counts per bin \\
$n\_cycles$ & int & Number of cycles \\
$bin\_duration$ & float & Bin duration (seconds) \\
$bin\_centers$ & array & Bin center times (seconds) \\
\hline
\multicolumn{3}{l}{\textbf{Turn Rate Formulas}} \\
\hline
$R$ & float & Turn rate (reorientations per minute, min$^{-1}$) \\
$N$ & int & Number of reorientation start events \\
$T$ & float & Time duration (seconds) \\
$R_{\text{bin}}$ & float & Bin turn rate (min$^{-1}$) \\
$N_{\text{bin},i}$ & int & Reorientation count in bin for cycle $i$ \\
$\Delta t_{\text{bin}}$ & float & Bin duration (seconds) \\
\hline
\multicolumn{3}{l}{\textbf{Angle Wrapping}} \\
\hline
$\theta_1$ & float & Initial angle (radians) \\
$\theta_2$ & float & Final angle (radians) \\
$\Delta\theta_{\text{raw}}$ & float & Raw angular difference (radians) \\
$\Delta\theta$ & float & Wrapped angular difference (radians, $[-\pi, +\pi]$) \\
$\pi$ & float & Half circle ($\pi = 3.14159...$ rad = $180°$) \\
$2\pi$ & float & Full circle ($2\pi = 6.28319...$ rad = $360°$) \\
\hline
\multicolumn{3}{l}{\textbf{Additional Formulas}} \\
\hline
$\bar{\theta}$ & float & Mean heading (circular mean, radians) \\
$L$ & float & Path length (cm or mm) \\
$\kappa$ & float & Curvature (rad/cm or rad/mm) \\
$s$ & float & Arc length (cm or mm) \\
$n$ & int & Number of samples \\
$i$ & int & Index or cycle number \\
\hline
\end{tabular}
\caption{Complete variable reference for the analysis pipeline}

\end{table}

\hypertarget{main-pipeline}{%
\section{Main Pipeline}\label{main-pipeline}}

\begin{algorithm}[H]
\caption{Main Analysis Pipeline}
\label{alg:main}
\begin{algorithmic}[1]
\REQUIRE H5 file path $h5\_file$
\ENSURE Events CSV, trajectories CSV, turn rate figures

\STATE Initialize progress monitor (PowerShell window)
\STATE Load H5 file: $h5\_data \leftarrow \text{load\_h5\_file}(h5\_file)$
\STATE Extract metadata: $fps \leftarrow h5\_data[\text{metadata}][\text{fps}]$
\STATE Extract stimulus timing: $stimulus\_df \leftarrow \text{extract\_stimulus\_timing}(h5\_data)$

\FOR{each track $t$ in $h5\_data[\text{tracks}]$}
    \STATE Process track $t$ (see Algorithm~\ref{alg:process-track})
\ENDFOR

\STATE Combine all event records: $events \leftarrow \text{concatenate}(all\_event\_records)$
\STATE Combine all trajectories: $trajectories \leftarrow \text{concatenate}(all\_trajectories)$
\STATE Save $events$ to CSV: $\text{experiment\_id}\_events.csv$
\STATE Save $trajectories$ to CSV: $\text{experiment\_id}\_trajectories.csv$

\STATE Calculate stimulus-locked turn rate (see Algorithm~\ref{alg:stimulus-locked})
\STATE Generate figures and save results
\end{algorithmic}
\end{algorithm}

\hypertarget{track-processing}{%
\section{Track Processing}\label{track-processing}}

\begin{algorithm}[H]
\caption{Process Single Track}
\label{alg:process-track}
\begin{algorithmic}[1]
\REQUIRE Track data $track\_data$, frame rate $fps$, stimulus data $stimulus\_df$
\ENSURE Event records $event\_records$, trajectory $traj\_df$

\STATE Extract trajectory features (see Algorithm~\ref{alg:extract-features})
\STATE Align trajectory with stimulus (see Algorithm~\ref{alg:align-stimulus})
\STATE Create event records (see Algorithm~\ref{alg:create-events})
\STATE Calculate per-cycle statistics (see Algorithm~\ref{alg:per-cycle-stats})

\RETURN $event\_records$, $traj\_df$
\end{algorithmic}
\end{algorithm}

\hypertarget{feature-extraction}{%
\section{Feature Extraction}\label{feature-extraction}}

\begin{algorithm}[H]
\caption{Extract Trajectory Features}
\label{alg:extract-features}
\begin{algorithmic}[1]
\REQUIRE Track data $track\_data$, frame rate $fps$
\ENSURE Trajectory DataFrame $df$ with features

\STATE Load raw data: $x, y, heading, speed, curvature \leftarrow track\_data$
\STATE Calculate heading change: $\Delta\theta[i] \leftarrow \text{wrap\_angle}(heading[i] - heading[i-1])$
\STATE Calculate acceleration: $accel[i] \leftarrow (speed[i] - speed[i-1]) / dt$

\STATE Run MAGAT segmentation (see Algorithm~\ref{alg:magat-segmentation})
\STATE Detect pauses (see Algorithm~\ref{alg:detect-pauses})
\STATE Extract turns from reorientations (see Algorithm~\ref{alg:extract-turns})

\STATE Create DataFrame $df$ with columns:
\STATE \quad $[time, x, y, heading, speed, curvature, \Delta\theta, accel,$
\STATE \quad $is\_run, is\_reorientation, is\_turn, is\_pause, ...]$

\RETURN $df$
\end{algorithmic}
\end{algorithm}

\hypertarget{magat-segmentation}{%
\section{MAGAT Segmentation}\label{magat-segmentation}}

\begin{algorithm}[H]
\caption{MAGAT Track Segmentation}
\label{alg:magat-segmentation}
\begin{algorithmic}[1]
\REQUIRE Trajectory data with $speed$, $curvature$, $body\_theta$, $vel\_dp$
\ENSURE Runs, head swings, reorientations

\STATE Define thresholds:
\STATE \quad $curv\_cut \leftarrow 0.4$ (curvature threshold)
\STATE \quad $theta\_cut \leftarrow \pi/2$ (body bend threshold)
\STATE \quad $stop\_speed\_cut \leftarrow 2.0$ mm/s
\STATE \quad $start\_speed\_cut \leftarrow 3.0$ mm/s

\STATE Find run end conditions:
\STATE \quad $high\_curv \leftarrow |curvature| > curv\_cut$
\STATE \quad $head\_swinging \leftarrow |body\_theta| > theta\_cut$
\STATE \quad $speed\_low \leftarrow speed < stop\_speed\_cut$
\STATE \quad $not\_a\_run \leftarrow high\_curv \vee head\_swinging \vee speed\_low$

\STATE Find run start conditions:
\STATE \quad $speed\_high \leftarrow speed \geq start\_speed\_cut$
\STATE \quad $head\_aligned \leftarrow vel\_dp \geq \cos(45°)$
\STATE \quad $is\_a\_run \leftarrow \neg not\_a\_run \wedge speed\_high \wedge head\_aligned$

\STATE Match run starts and ends to create run intervals
\STATE Apply quality filters:
\STATE \quad Filter runs with duration $< minRunTime$ (default: 2.5s)
\STATE \quad Filter runs with path length $< minRunLength$ (optional)
\STATE \quad Filter runs with avg speed $< minRunSpeed$ (optional)

\STATE Detect head swings:
\STATE \quad $head\_swing\_start \leftarrow |body\_theta| > 20°$ AND $\neg is\_run$
\STATE \quad $head\_swing\_stop \leftarrow |body\_theta| < 10°$ OR sign change
\STATE \quad Match starts and stops to create head swing intervals

\STATE Detect reorientations:
\STATE \quad Reorientations = gaps between runs (periods where $\neg is\_run$)

\RETURN $runs$, $head\_swings$, $reorientations$
\end{algorithmic}
\end{algorithm}

\hypertarget{pause-detection}{%
\section{Pause Detection}\label{pause-detection}}

\begin{algorithm}[H]
\caption{Detect Pauses}
\label{alg:detect-pauses}
\begin{algorithmic}[1]
\REQUIRE Speed array $speed$, time array $time$
\ENSURE Pause boolean array $is\_pause$

\STATE Define thresholds:
\STATE \quad $speed\_threshold \leftarrow 0.001$ (mm/s or cm/s)
\STATE \quad $min\_duration \leftarrow 0.2$ seconds

\STATE $is\_pause \leftarrow speed < speed\_threshold$

\STATE Find pause start/end events:
\FOR{$i = 1$ to $len(speed)$}
    \IF{$is\_pause[i]$ AND $\neg is\_pause[i-1]$}
        \STATE $pause\_starts.\text{append}(i)$
    \ELSIF{$\neg is\_pause[i]$ AND $is\_pause[i-1]$}
        \STATE $pause\_ends.\text{append}(i)$
    \ENDIF
\ENDFOR

\STATE Filter by minimum duration:
\FOR{each pause $(start, end)$}
    \STATE $duration \leftarrow time[end] - time[start]$
    \IF{$duration \geq min\_duration$}
        \STATE Mark frames $[start:end]$ as valid pause
    \ENDIF
\ENDFOR

\RETURN $is\_pause$
\end{algorithmic}
\end{algorithm}

\hypertarget{turn-extraction}{%
\section{Turn Extraction}\label{turn-extraction}}

\begin{algorithm}[H]
\caption{Extract Turns from Reorientations}
\label{alg:extract-turns}
\begin{algorithmic}[1]
\REQUIRE Reorientations $reorientations$, pause array $is\_pause$
\ENSURE Turns array $turns$, boolean array $is\_turn$

\STATE Initialize: $turns \leftarrow []$, $is\_turn \leftarrow \text{zeros}(n\_frames)$

\FOR{each reorientation $(reo\_start, reo\_end)$ in $reorientations$}
    \STATE Extract pause frames in reorientation:
    \STATE \quad $reo\_pause\_frames \leftarrow is\_pause[reo\_start:reo\_end+1]$
    
    \IF{$\exists$ pause in $reo\_pause\_frames$}
        \STATE This reorientation contains a pause $\Rightarrow$ it's a turn
        \STATE $turns.\text{append}((reo\_start, reo\_end))$
        \STATE $is\_turn[reo\_start:reo\_end+1] \leftarrow \text{True}$
    \ENDIF
\ENDFOR

\RETURN $turns$, $is\_turn$
\end{algorithmic}
\end{algorithm}

\hypertarget{stimulus-alignment}{%
\section{Stimulus Alignment}\label{stimulus-alignment}}

\begin{algorithm}[H]
\caption{Align Trajectory with Stimulus}
\label{alg:align-stimulus}
\begin{algorithmic}[1]
\REQUIRE Trajectory $traj\_df$, stimulus data $stimulus\_df$
\ENSURE Aligned trajectory $aligned\_df$

\STATE For each row in $traj\_df$:
\FOR{each time $t$ in $traj\_df[\text{time}]$}
    \STATE Find nearest stimulus time: $stim\_idx \leftarrow \text{argmin}(|stimulus\_df[\text{time}] - t|)$
    \STATE Interpolate LED values: $led1Val \leftarrow \text{interpolate}(stimulus\_df[\text{led1Val}], t)$
    \STATE Calculate period-relative timing:
    \STATE \quad $led12Val\_ton \leftarrow$ time within LED ON period
    \STATE \quad $led12Val\_toff \leftarrow$ time within LED OFF period
    \STATE Add stimulus flags:
    \STATE \quad $stimulus\_on \leftarrow led1Val > threshold$
    \STATE \quad $stimulus\_onset \leftarrow$ transition from OFF to ON
\ENDFOR

\RETURN $aligned\_df$
\end{algorithmic}
\end{algorithm}

\hypertarget{event-record-creation}{%
\section{Event Record Creation}\label{event-record-creation}}

\begin{algorithm}[H]
\caption{Create Event Records}
\label{alg:create-events}
\begin{algorithmic}[1]
\REQUIRE Aligned trajectory $aligned\_df$, track ID $track\_id$, experiment ID $exp\_id$
\ENSURE Event records DataFrame $event\_records$

\STATE Define bin size: $bin\_size \leftarrow 0.05$ seconds

\STATE Sort trajectory by time: $aligned\_df \leftarrow \text{sort}(aligned\_df, \text{by}=time)$

\STATE Detect START events (False $\rightarrow$ True transitions):
\STATE \quad $is\_reo\_padded \leftarrow [\text{False}] + is\_reorientation$
\STATE \quad $reo\_start\_mask \leftarrow \neg is\_reo\_padded[:-1] \wedge is\_reo\_padded[1:]$
\STATE \quad $is\_reorientation\_start[reo\_start\_mask] \leftarrow \text{True}$

\STATE Repeat for $is\_turn\_start$ and $is\_pause\_start$

\STATE Bin trajectory into $bin\_size$ bins:
\FOR{each bin $b$}
    \STATE $bin\_start \leftarrow b \times bin\_size$
    \STATE $bin\_end \leftarrow (b+1) \times bin\_size$
    \STATE $bin\_data \leftarrow aligned\_df[\text{where}(time \in [bin\_start, bin\_end))]$
    
    \STATE Aggregate:
    \STATE \quad $time \leftarrow \text{first}(bin\_data[\text{time}])$
    \STATE \quad $is\_reorientation\_start \leftarrow \text{any}(bin\_data[\text{is\_reorientation\_start}])$
    \STATE \quad $is\_turn\_start \leftarrow \text{any}(bin\_data[\text{is\_turn\_start}])$
    \STATE \quad $led1Val \leftarrow \text{mean}(bin\_data[\text{led1Val}])$
    \STATE \quad ... (other aggregations)
\ENDFOR

\RETURN $event\_records$
\end{algorithmic}
\end{algorithm}

\hypertarget{per-cycle-statistics}{%
\section{Per-Cycle Statistics}\label{per-cycle-statistics}}

\begin{algorithm}[H]
\caption{Calculate Per-Cycle Statistics}
\label{alg:per-cycle-stats}
\begin{algorithmic}[1]
\REQUIRE H5 file $h5\_file$, aligned trajectory $aligned\_df$
\ENSURE Console output with per-cycle statistics

\STATE Extract cycles from H5:
\STATE \quad $cycles \leftarrow \text{extract\_cycles\_from\_h5}(h5\_file)$

\STATE Extract reorientation start times:
\STATE \quad $reo\_starts \leftarrow aligned\_df[\text{where}(is\_reorientation\_start == \text{True})]$

\STATE Extract turn start times:
\STATE \quad $turn\_starts \leftarrow aligned\_df[\text{where}(is\_turn\_start == \text{True})]$

\FOR{each cycle $c$ in $cycles$}
    \STATE $cycle\_start \leftarrow c[\text{cycle\_start\_time}]$
    \STATE $cycle\_end \leftarrow c[\text{cycle\_end\_time}]$
    \STATE $onset\_time \leftarrow c[\text{onset\_time}]$
    
    \STATE Count reorientations in cycle:
    \STATE \quad $cycle\_reos \leftarrow reo\_starts[\text{where}(time \in [cycle\_start, cycle\_end])]$
    \STATE \quad $n\_cycle\_reos \leftarrow \text{len}(cycle\_reos)$
    
    \STATE Count turns in cycle:
    \STATE \quad $cycle\_turns \leftarrow turn\_starts[\text{where}(time \in [cycle\_start, cycle\_end])]$
    \STATE \quad $n\_cycle\_turns \leftarrow \text{len}(cycle\_turns)$
    
    \STATE Calculate cycle turn rate:
    \STATE \quad $cycle\_duration\_min \leftarrow (cycle\_end - cycle\_start) / 60.0$
    \STATE \quad $turn\_rate \leftarrow n\_cycle\_reos / cycle\_duration\_min$
    
    \STATE Print cycle summary header
    
    \STATE Calculate per-bin statistics:
    \STATE \quad $bin\_size \leftarrow 0.5$ seconds
    \STATE \quad $n\_bins \leftarrow \lceil (cycle\_end - cycle\_start) / bin\_size \rceil$
    
    \FOR{each bin $b = 0$ to $n\_bins-1$}
        \STATE $bin\_start \leftarrow cycle\_start + b \times bin\_size$
        \STATE $bin\_end \leftarrow \min(cycle\_start + (b+1) \times bin\_size, cycle\_end)$
        
        \STATE Count reorientations in bin:
        \STATE \quad $bin\_reos \leftarrow reo\_starts[\text{where}(time \in [bin\_start, bin\_end))]$
        \STATE \quad $n\_bin\_reos \leftarrow \text{len}(bin\_reos)$
        
        \STATE Count turns in bin:
        \STATE \quad $n\_bin\_turns \leftarrow \text{len}(turn\_starts[\text{where}(time \in [bin\_start, bin\_end))])$
        
        \STATE Calculate bin turn rate:
        \STATE \quad $bin\_duration\_min \leftarrow (bin\_end - bin\_start) / 60.0$
        \STATE \quad $bin\_rate \leftarrow n\_bin\_reos / bin\_duration\_min$
        
        \STATE Calculate time relative to onset:
        \STATE \quad $time\_rel\_onset \leftarrow (bin\_start + bin\_end)/2 - onset\_time$
        
        \STATE Print bin statistics: $b$, $time\_rel\_onset$, $n\_bin\_reos$, $n\_bin\_turns$, $bin\_rate$
    \ENDFOR
\ENDFOR

\STATE Print track totals
\end{algorithmic}
\end{algorithm}

\hypertarget{stimulus-locked-turn-rate-calculation}{%
\section{Stimulus-Locked Turn Rate
Calculation}\label{stimulus-locked-turn-rate-calculation}}

\begin{algorithm}[H]
\caption{Calculate Stimulus-Locked Turn Rate}
\label{alg:stimulus-locked}
\begin{algorithmic}[1]
\REQUIRE Events CSV $events\_df$, trajectories CSV $trajectories\_df$, H5 file $h5\_file$
\ENSURE Turn rate PSTH (Peri-Stimulus Time Histogram)

\STATE Extract cycles: $cycles \leftarrow \text{extract\_cycles\_from\_h5}(h5\_file)$

\STATE Extract reorientation start times:
\IF{$is\_reorientation\_start$ column exists in $events\_df$}
    \STATE Use $events\_df[\text{is\_reorientation\_start}]$ directly
\ELSIF{$trajectories\_df$ provided AND $is\_reorientation$ exists}
    \STATE Detect False $\rightarrow$ True transitions in $trajectories\_df$
    \STATE Extract start times for each track
\ELSE
    \STATE Fallback to $events\_df[\text{is\_reorientation}]$ (less accurate)
\ENDIF

\STATE Define analysis bins: $bin\_edges \leftarrow [-10, -9.5, ..., 0, 0.5, ..., pulse\_duration]$

\FOR{each cycle $c$ in $cycles$}
    \STATE $onset\_time \leftarrow c[\text{onset\_time}]$
    \STATE $cycle\_start \leftarrow c[\text{cycle\_start\_time}]$
    \STATE $cycle\_end \leftarrow c[\text{cycle\_end\_time}]$
    
    \FOR{each reorientation start time $t\_reo$}
        \STATE Calculate time relative to onset: $t\_rel \leftarrow t\_reo - onset\_time$
        \IF{$t\_rel \in [cycle\_start - onset\_time, cycle\_end - onset\_time]$}
            \STATE Find bin: $bin\_idx \leftarrow \text{find\_bin}(t\_rel, bin\_edges)$
            \STATE Increment count: $counts[bin\_idx] \leftarrow counts[bin\_idx] + 1$
        \ENDIF
    \ENDFOR
\ENDFOR

\STATE Calculate turn rate per bin:
\FOR{each bin $b$}
    \STATE $bin\_duration \leftarrow bin\_edges[b+1] - bin\_edges[b]$ (seconds)
    \STATE $bin\_duration\_min \leftarrow bin\_duration / 60.0$
    \STATE $turn\_rate[b] \leftarrow counts[b] / (n\_cycles \times bin\_duration\_min)$
\ENDFOR

\RETURN $turn\_rate$, $bin\_centers$
\end{algorithmic}
\end{algorithm}

\hypertarget{key-formulas}{%
\section{Key Formulas}\label{key-formulas}}

\hypertarget{turn-rate-calculation}{%
\subsection{Turn Rate Calculation}\label{turn-rate-calculation}}

\hypertarget{formula}{%
\subsubsection{Formula}\label{formula}}

The turn rate (reorientations per minute) is calculated as:

\[R = \frac{N}{T \times 60}\]

\hypertarget{term-breakdown}{%
\subsubsection{Term Breakdown}\label{term-breakdown}}

\begin{table}[H]
\centering
\small
\begin{tabular}{p{3cm}p{11cm}}
\hline
\textbf{Term} & \textbf{Description} \\
\hline
$R$ (Turn Rate) & \textbf{Definition}: Average number of reorientation start events per minute \\
 & \textbf{Units}: min$^{-1}$ (reorientations per minute) \\
 & \textbf{Typical Range}: 0.1 - 10 min$^{-1}$ for larval behavior \\
 & \textbf{Resolution}: Limited by bin size (default: 0.05s bins) and frame rate (typically 10 Hz) \\
 & \textbf{Interpretation}: Higher values indicate more frequent reorientation behavior \\
\hline
$N$ (Reorientation Starts) & \textbf{Definition}: Count of reorientation START events (False $\rightarrow$ True transitions) \\
 & \textbf{Units}: dimensionless (count) \\
 & \textbf{Critical Note}: Only START events are counted, NOT all frames during reorientations \\
 & \textbf{Detection}: Pad array: $[\text{False}] + is\_reorientation$, detect transitions, count \\
 & \textbf{Resolution}: Frame-level accuracy (typically 0.1s at 10 Hz) \\
 & \textbf{Example}: If reorientation lasts 2s (20 frames), only 1 event counted (the start) \\
\hline
$T$ (Time Duration) & \textbf{Definition}: Total time period over which reorientations are counted \\
 & \textbf{Units}: seconds \\
 & \textbf{Typical Values}: Per track: 5-30 min; Per cycle: 10-40s; Per bin: 0.5s \\
 & \textbf{Resolution}: Limited by frame rate (typically 0.1s at 10 Hz) \\
 & \textbf{Calculation}: $T = t_{\text{end}} - t_{\text{start}}$ where $t$ is in seconds \\
\hline
60 (Conversion) & \textbf{Purpose}: Convert from events per second to events per minute \\
 & \textbf{Derivation}: 1 minute = 60 seconds, so $\frac{\text{events}}{\text{second}} \times 60 = \frac{\text{events}}{\text{minute}}$ \\
\hline
\end{tabular}
\caption{Term breakdown for turn rate calculation}

\end{table}

\hypertarget{formula-derivation}{%
\subsubsection{Formula Derivation}\label{formula-derivation}}

Starting from the definition of rate:

\[R = \frac{\text{number of events}}{\text{time period}}\]

Substituting our terms:

\[R = \frac{N}{T}\]

To convert to per-minute units:

\[R = \frac{N}{T} \times \frac{60 \text{ seconds}}{1 \text{ minute}} = \frac{N}{T \times 60}\]

\hypertarget{example-calculation}{%
\subsubsection{Example Calculation}\label{example-calculation}}

\begin{table}[H]
\centering
\small
\begin{tabular}{ll}
\hline
\textbf{Given} & \textbf{Value} \\
\hline
$N$ (reorientation starts) & 25 events \\
$T$ (time duration) & 5.0 minutes = 300 seconds \\
\hline
\textbf{Calculation} & \\
\hline
If $T$ in minutes: & $R = \frac{N}{T} = \frac{25}{5} = 5.0 \text{ min}^{-1}$ \\
If $T$ in seconds: & $R = \frac{N}{T} \times 60 = \frac{25}{300} \times 60 = 5.0 \text{ min}^{-1}$ \\
\hline
\end{tabular}
\caption{Example turn rate calculation}

\end{table}

\hypertarget{resolution-and-precision}{%
\subsubsection{Resolution and
Precision}\label{resolution-and-precision}}

\begin{table}[H]
\centering
\small
\begin{tabular}{p{5cm}p{9cm}}
\hline
\textbf{Aspect} & \textbf{Value} \\
\hline
Temporal Resolution & Limited by frame rate (10 Hz = 0.1s) \\
Event Detection Resolution & Frame-level (0.1s at 10 Hz) \\
Statistical Precision & Improves with longer observation time $T$ \\
Minimum Detectable Rate & $\frac{1}{T \times 60}$ min$^{-1}$ (one event in duration $T$) \\
\hline
\end{tabular}
\caption{Resolution and precision for turn rate calculation}

\end{table}

\begin{center}\rule{0.5\linewidth}{0.5pt}\end{center}

\hypertarget{stimulus-locked-turn-rate}{%
\subsection{Stimulus-Locked Turn Rate}\label{stimulus-locked-turn-rate}}

\hypertarget{formula-1}{%
\subsubsection{Formula}\label{formula-1}}

For stimulus-locked analysis, the rate is calculated per bin:

\[R_{\text{bin}} = \frac{\sum_{i=1}^{n_{\text{cycles}}} N_{\text{bin},i}}{n_{\text{cycles}} \times \Delta t_{\text{bin}} \times 60}\]

\hypertarget{term-breakdown-1}{%
\subsubsection{Term Breakdown}\label{term-breakdown-1}}

\begin{table}[H]
\centering
\footnotesize
\begin{tabular}{p{3cm}p{11cm}}
\hline
\textbf{Term} & \textbf{Description} \\
\hline
$R_{\text{bin}}$ & \textbf{Definition}: Average turn rate for specific time bin relative to stimulus onset \\
 & \textbf{Units}: min$^{-1}$ (reorientations per minute) \\
 & \textbf{Range}: 0 - 20 min$^{-1}$ (can spike during stimulus) \\
 & \textbf{Resolution}: Limited by bin size (0.5s per-cycle, 0.05s event records) \\
 & \textbf{Use}: Creates Peri-Stimulus Time Histogram (PSTH) \\
\hline
$N_{\text{bin},i}$ & \textbf{Definition}: Count of reorientation starts in bin for cycle $i$ \\
 & \textbf{Units}: dimensionless (count) \\
 & \textbf{Detection}: For cycle $i$ with onset $t_{\text{onset},i}$, count events in \\
 & \quad $[t_{\text{onset},i} + t_{\text{bin\_start}}, t_{\text{onset},i} + t_{\text{bin\_end}})$ \\
 & \textbf{Range}: 0, 1, 2, ... (typically 0-5 per bin per cycle) \\
\hline
$n_{\text{cycles}}$ & \textbf{Definition}: Total number of stimulus cycles \\
 & \textbf{Units}: dimensionless (count) \\
 & \textbf{Typical}: 10-50 cycles per experiment \\
 & \textbf{Impact}: More cycles $\Rightarrow$ better statistical precision \\
\hline
$\Delta t_{\text{bin}}$ & \textbf{Definition}: Time width of analysis bin \\
 & \textbf{Units}: seconds \\
 & \textbf{Typical}: 0.5s (per-cycle), 0.05s (events), 0.5-1.0s (PSTH) \\
 & \textbf{Trade-off}: Smaller bins = higher resolution but noisier \\
\hline
$\sum N_{\text{bin},i}$ & \textbf{Definition}: Sum of reorientation counts across all cycles \\
 & \textbf{Units}: dimensionless (count) \\
 & \textbf{Range}: 0 - 100 (depends on $n_{\text{cycles}}$ and behavior) \\
\hline
60 (Conversion) & \textbf{Purpose}: Convert events/second to events/minute \\
\hline
\end{tabular}
\caption{Term breakdown for stimulus-locked turn rate}

\end{table}

\hypertarget{formula-derivation-1}{%
\subsubsection{Formula Derivation}\label{formula-derivation-1}}

Starting from the basic turn rate formula applied to a single bin:

\[R_{\text{bin},i} = \frac{N_{\text{bin},i}}{\Delta t_{\text{bin}} \times 60}\]

This gives the rate for cycle \(i\) in this bin. To average across all
cycles:

\[R_{\text{bin}} = \frac{1}{n_{\text{cycles}}} \sum_{i=1}^{n_{\text{cycles}}} R_{\text{bin},i}\]

Substituting:

\[R_{\text{bin}} = \frac{1}{n_{\text{cycles}}} \sum_{i=1}^{n_{\text{cycles}}} \frac{N_{\text{bin},i}}{\Delta t_{\text{bin}} \times 60}\]

Rearranging:

\[R_{\text{bin}} = \frac{\sum_{i=1}^{n_{\text{cycles}}} N_{\text{bin},i}}{n_{\text{cycles}} \times \Delta t_{\text{bin}} \times 60}\]

\hypertarget{example-calculation-1}{%
\subsubsection{Example Calculation}\label{example-calculation-1}}

\begin{table}[H]
\centering
\small
\begin{tabular}{ll}
\hline
\textbf{Given} & \textbf{Value} \\
\hline
$n_{\text{cycles}}$ & 20 cycles \\
$\Delta t_{\text{bin}}$ & 0.5 seconds \\
Bin counts $N_{\text{bin},i}$ & $[0, 1, 2, 0, 1, 0, 3, 1, 0, 2, 1, 0, 1, 2, 0, 1, 1, 0, 2, 1]$ \\
$\sum N_{\text{bin},i}$ & 21 reorientations \\
\hline
\textbf{Calculation} & \\
\hline
Direct formula: & $R_{\text{bin}} = \frac{21}{20 \times 0.5 \times 60} = \frac{21}{600} = 0.035 \text{ min}^{-1}$ \\
\hline
\textbf{Alternative (intuitive)} & \\
\hline
Average per cycle: & $\bar{N} = \frac{21}{20} = 1.05$ reorientations/cycle \\
Rate per cycle: & $R_{\text{cycle}} = \frac{1.05}{0.5 \times 60} = 0.035 \text{ min}^{-1}$ \\
\hline
\end{tabular}
\caption{Example stimulus-locked turn rate calculation}

\end{table}

\hypertarget{resolution-and-precision-1}{%
\subsubsection{Resolution and
Precision}\label{resolution-and-precision-1}}

\begin{table}[H]
\centering
\small
\begin{tabular}{p{5cm}p{9cm}}
\hline
\textbf{Aspect} & \textbf{Value} \\
\hline
Temporal Resolution & Limited by bin size $\Delta t_{\text{bin}}$ and frame rate \\
Spatial Resolution & Limited by frame rate (0.1s at 10 Hz) \\
Statistical Precision & Standard error: $\text{SE} = \frac{\sqrt{\sum N_{\text{bin},i}}}{n_{\text{cycles}} \times \Delta t_{\text{bin}} \times 60}$ \\
 & Improves with $\sqrt{n_{\text{cycles}}}$ (more cycles = better precision) \\
Minimum Detectable Rate & $\frac{1}{n_{\text{cycles}} \times \Delta t_{\text{bin}} \times 60}$ min$^{-1}$ \\
\hline
\end{tabular}
\caption{Resolution and precision for stimulus-locked turn rate}

\end{table}

\hypertarget{bin-alignment}{%
\subsubsection{Bin Alignment}\label{bin-alignment}}

\begin{table}[H]
\centering
\small
\begin{tabular}{p{4cm}p{10cm}}
\hline
\textbf{Bin Type} & \textbf{Time Range} \\
\hline
Baseline bins & $t \in [-10, 0)$ seconds (before stimulus) \\
Stimulus bins & $t \in [0, \text{pulse\_duration}]$ seconds (during stimulus) \\
Example edges & $[-10, -9.5, -9.0, ..., -0.5, 0, 0.5, ..., 30.0]$ seconds \\
\hline
\end{tabular}
\caption{Bin alignment relative to stimulus onset (t=0)}

\end{table}

\begin{center}\rule{0.5\linewidth}{0.5pt}\end{center}

\hypertarget{angle-wrapping}{%
\subsection{Angle Wrapping}\label{angle-wrapping}}

\hypertarget{formula-2}{%
\subsubsection{Formula}\label{formula-2}}

Angular differences are wrapped to \([-\pi, +\pi]\):

\[\Delta\theta = \begin{cases}
\theta_2 - \theta_1 + 2\pi & \text{if } \Delta\theta_{\text{raw}} < -\pi \\
\theta_2 - \theta_1 - 2\pi & \text{if } \Delta\theta_{\text{raw}} > +\pi \\
\theta_2 - \theta_1 & \text{otherwise}
\end{cases}\]

where \(\Delta\theta_{\text{raw}} = \theta_2 - \theta_1\) (unwrapped
difference)

\hypertarget{term-breakdown-2}{%
\subsubsection{Term Breakdown}\label{term-breakdown-2}}

\begin{table}[H]
\centering
\footnotesize
\begin{tabular}{p{3cm}p{11cm}}
\hline
\textbf{Term} & \textbf{Description} \\
\hline
$\theta_1$ (Initial) & \textbf{Definition}: Heading angle at start of measurement period \\
 & \textbf{Units}: radians \\
 & \textbf{Range}: $[0, 2\pi)$ or $(-\pi, +\pi]$ depending on convention \\
 & \textbf{Resolution}: Limited by tracking precision (typically 0.01-0.1 radians) \\
\hline
$\theta_2$ (Final) & \textbf{Definition}: Heading angle at end of measurement period \\
 & \textbf{Units}: radians \\
 & \textbf{Range}: $[0, 2\pi)$ or $(-\pi, +\pi]$ depending on convention \\
 & \textbf{Resolution}: Same as $\theta_1$ \\
\hline
$\Delta\theta_{\text{raw}}$ & \textbf{Definition}: Unwrapped angular difference: $\theta_2 - \theta_1$ \\
 & \textbf{Units}: radians \\
 & \textbf{Range}: $(-\infty, +\infty)$ (can exceed $[-\pi, +\pi]$) \\
 & \textbf{Problem}: Exceeds $[-\pi, +\pi]$ when angles wrap around $2\pi$ boundary \\
 & \textbf{Example}: $\theta_1 = 350°$, $\theta_2 = 10°$ $\Rightarrow$ $\Delta\theta_{\text{raw}} = -6.0$ rad (wrong) \\
\hline
$\Delta\theta$ (Wrapped) & \textbf{Definition}: Angular difference wrapped to $[-\pi, +\pi]$ \\
 & \textbf{Units}: radians \\
 & \textbf{Range}: $[-\pi, +\pi]$ (always) \\
 & \textbf{Interpretation}: Positive = counter-clockwise (left), Negative = clockwise (right) \\
 & \textbf{Resolution}: Same as angle measurements (0.01-0.1 radians) \\
\hline
$2\pi$ (Full Circle) & \textbf{Value}: $2\pi = 6.283185307...$ rad = $360°$ \\
 & \textbf{Purpose}: Used to wrap angles back into $[-\pi, +\pi]$ range \\
\hline
$\pi$ (Half Circle) & \textbf{Value}: $\pi = 3.141592654...$ rad = $180°$ \\
 & \textbf{Purpose}: Boundary of wrapped range $[-\pi, +\pi]$ \\
\hline
\end{tabular}
\caption{Term breakdown for angle wrapping}

\end{table}

\hypertarget{formula-derivation-2}{%
\subsubsection{Formula Derivation}\label{formula-derivation-2}}

The wrapping formula ensures the shortest angular path:

\begin{enumerate}
\def\labelenumi{\arabic{enumi}.}
\item
  \textbf{Calculate raw difference}:
  \(\Delta\theta_{\text{raw}} = \theta_2 - \theta_1\)
\item
  \textbf{Check if wrapping needed}:

  \begin{itemize}
  \tightlist
  \item
    If \(\Delta\theta_{\text{raw}} < -\pi\): Add \(2\pi\) to get
    equivalent angle in \([-\pi, +\pi]\)
  \item
    If \(\Delta\theta_{\text{raw}} > +\pi\): Subtract \(2\pi\) to get
    equivalent angle in \([-\pi, +\pi]\)
  \item
    Otherwise: Use raw difference
  \end{itemize}
\item
  \textbf{Mathematical justification}: Angles are periodic modulo
  \(2\pi\), so:

  \begin{itemize}
  \tightlist
  \item
    \(\theta + 2\pi \equiv \theta\) (same direction)
  \item
    Adding/subtracting \(2\pi\) doesn't change the actual direction
  \end{itemize}
\end{enumerate}

\hypertarget{example-calculations}{%
\subsubsection{Example Calculations}\label{example-calculations}}

\begin{table}[H]
\centering
\footnotesize
\begin{tabular}{p{2.5cm}p{5cm}p{6.5cm}}
\hline
\textbf{Example} & \textbf{Input} & \textbf{Result} \\
\hline
1: No wrapping & $\theta_1 = 45°$, $\theta_2 = 90°$ & $\Delta\theta_{\text{raw}} = 0.786$ rad $< \pi$ \\
 & & $\Delta\theta = 0.786$ rad = $45°$ (counter-clockwise) \\
\hline
2: Wrap negative & $\theta_1 = 350°$, $\theta_2 = 10°$ & $\Delta\theta_{\text{raw}} = -5.933$ rad $< -\pi$ \\
 & & Add $2\pi$: $\Delta\theta = +0.350$ rad = $+20°$ \\
 & & Shortest path: $+20°$ (not $-340°$) \\
\hline
3: Wrap positive & $\theta_1 = 10°$, $\theta_2 = 350°$ & $\Delta\theta_{\text{raw}} = +5.933$ rad $> +\pi$ \\
 & & Subtract $2\pi$: $\Delta\theta = -0.350$ rad = $-20°$ \\
 & & Shortest path: $-20°$ (not $+340°$) \\
\hline
4: Large turn & $\theta_1 = 0°$, $\theta_2 = 200°$ & $\Delta\theta_{\text{raw}} = +3.490$ rad $> +\pi$ \\
 & & Subtract $2\pi$: $\Delta\theta = -2.793$ rad = $-160°$ \\
 & & $200°$ turn $\equiv$ $-160°$ turn (shorter path) \\
\hline
\end{tabular}
\caption{Example calculations for angle wrapping}

\end{table}

\hypertarget{implementation}{%
\subsubsection{Implementation}\label{implementation}}

In Python/NumPy:

\begin{Shaded}
\begin{Highlighting}[]
\ImportTok{import}\NormalTok{ numpy }\ImportTok{as}\NormalTok{ np}

\KeywordTok{def}\NormalTok{ wrap\_angle\_diff(theta1, theta2):}
    \CommentTok{"""}
\CommentTok{    Calculate angular difference with proper wrapping.}
\CommentTok{    }
\CommentTok{    Parameters:}
\CommentTok{    {-}{-}{-}{-}{-}{-}{-}{-}{-}{-}{-}}
\CommentTok{    theta1 : array{-}like}
\CommentTok{        Initial angles in radians}
\CommentTok{    theta2 : array{-}like}
\CommentTok{        Final angles in radians}
\CommentTok{    }
\CommentTok{    Returns:}
\CommentTok{    {-}{-}{-}{-}{-}{-}{-}{-}}
\CommentTok{    delta\_theta : ndarray}
\CommentTok{        Wrapped angular difference in [{-}π, +π]}
\CommentTok{    """}
\NormalTok{    delta\_theta }\OperatorTok{=}\NormalTok{ theta2 }\OperatorTok{{-}}\NormalTok{ theta1}
    
    \CommentTok{\# Wrap to [{-}π, +π]}
\NormalTok{    delta\_theta }\OperatorTok{=}\NormalTok{ np.where(delta\_theta }\OperatorTok{\textless{}} \OperatorTok{{-}}\NormalTok{np.pi, }
\NormalTok{                           delta\_theta }\OperatorTok{+} \DecValTok{2}\OperatorTok{*}\NormalTok{np.pi, }
\NormalTok{                           delta\_theta)}
\NormalTok{    delta\_theta }\OperatorTok{=}\NormalTok{ np.where(delta\_theta }\OperatorTok{\textgreater{}}\NormalTok{ np.pi, }
\NormalTok{                           delta\_theta }\OperatorTok{{-}} \DecValTok{2}\OperatorTok{*}\NormalTok{np.pi, }
\NormalTok{                           delta\_theta)}
    
    \ControlFlowTok{return}\NormalTok{ delta\_theta}
\end{Highlighting}
\end{Shaded}

\hypertarget{resolution-and-precision-2}{%
\subsubsection{Resolution and
Precision}\label{resolution-and-precision-2}}

\begin{table}[H]
\centering
\small
\begin{tabular}{p{5cm}p{9cm}}
\hline
\textbf{Aspect} & \textbf{Value} \\
\hline
Angular Resolution & Limited by tracking precision \\
 & Typical: 0.01-0.1 radians ($0.6°$ - $6°$) \\
 & High-quality: 0.001 radians ($0.06°$) \\
Wrapping Precision & Exact (no approximation) \\
Edge Cases & Exactly $\pm\pi$: Both directions equally valid (choose $+\pi$) \\
 & Very small differences: Wrapping doesn't affect (difference $< \pi$) \\
\hline
\end{tabular}
\caption{Resolution and precision for angle wrapping}

\end{table}

\hypertarget{use-cases}{%
\subsubsection{Use Cases}\label{use-cases}}

\begin{table}[H]
\centering
\small
\begin{tabular}{p{4cm}p{10cm}}
\hline
\textbf{Use Case} & \textbf{Formula} \\
\hline
Heading Change & $\Delta\theta = \text{wrap\_angle\_diff}(heading[i-1], heading[i])$ \\
Turn Magnitude & $|\Delta\theta|$ gives turn size (independent of direction) \\
Turn Direction & $\text{sign}(\Delta\theta)$ gives direction (positive = left, negative = right) \\
Run Drift & $\Delta\theta = \text{wrap\_angle\_diff}(run\_start\_theta, run\_end\_theta)$ \\
Reorientation Magnitude & $\Delta\theta = \text{wrap\_angle\_diff}(reo\_prev\_dir, reo\_next\_dir)$ \\
\hline
\end{tabular}
\caption{Use cases for angle wrapping}

\end{table}

\begin{center}\rule{0.5\linewidth}{0.5pt}\end{center}

\hypertarget{additional-formulas}{%
\subsection{Additional Formulas}\label{additional-formulas}}

\hypertarget{mean-heading-circular-mean}{%
\subsubsection{Mean Heading (Circular
Mean)}\label{mean-heading-circular-mean}}

For averaging angles (e.g., run mean heading):

\[\bar{\theta} = \text{atan2}\left(\frac{1}{n}\sum_{i=1}^{n} \sin(\theta_i), \frac{1}{n}\sum_{i=1}^{n} \cos(\theta_i)\right)\]

\begin{table}[H]
\centering
\small
\begin{tabular}{p{3cm}p{11cm}}
\hline
\textbf{Aspect} & \textbf{Description} \\
\hline
Purpose & Average angles accounting for circularity ($0° \equiv 360°$) \\
Method & Convert to Cartesian coordinates, average, convert back \\
Why needed & Simple arithmetic mean fails (e.g., mean of $350°$ and $10°$ should be $0°$, not $180°$) \\
\hline
\end{tabular}
\caption{Circular mean heading breakdown}

\end{table}

\hypertarget{path-length}{%
\subsubsection{Path Length}\label{path-length}}

Cumulative distance along trajectory:

\[L = \sum_{i=1}^{n-1} \sqrt{(x_{i+1} - x_i)^2 + (y_{i+1} - y_i)^2}\]

\begin{table}[H]
\centering
\small
\begin{tabular}{p{3cm}p{11cm}}
\hline
\textbf{Aspect} & \textbf{Description} \\
\hline
Purpose & Total distance traveled (not Euclidean displacement) \\
Units & cm or mm (same as position units) \\
Resolution & Frame-level (typically 0.01-0.1 cm per frame) \\
Use & Run length, total path length \\
\hline
\end{tabular}
\caption{Path length breakdown}

\end{table}

\hypertarget{curvature}{%
\subsubsection{Curvature}\label{curvature}}

Rate of change of heading:

\[\kappa = \frac{d\theta}{ds} = \frac{\Delta\theta}{\Delta s}\]

where \(s\) is arc length.

\begin{table}[H]
\centering
\small
\begin{tabular}{p{3cm}p{11cm}}
\hline
\textbf{Aspect} & \textbf{Description} \\
\hline
Purpose & Measure of how sharply trajectory curves \\
Units & rad/cm or rad/mm \\
Typical Range & 0 (straight) to 10 rad/cm (very curved) \\
Use & Run detection (high curvature ends runs) \\
\hline
\end{tabular}
\caption{Curvature breakdown}

\end{table}

\hypertarget{critical-implementation-notes}{%
\section{Critical Implementation
Notes}\label{critical-implementation-notes}}

\hypertarget{start-event-detection}{%
\subsection{START Event Detection}\label{start-event-detection}}

\textbf{CRITICAL}: Only count START events (False \(\rightarrow\) True
transitions), not all frames during events.

\begin{Shaded}
\begin{Highlighting}[]
\CommentTok{\# Correct approach:}
\NormalTok{is\_reo\_padded }\OperatorTok{=}\NormalTok{ [}\VariableTok{False}\NormalTok{] }\OperatorTok{+}\NormalTok{ is\_reorientation}
\NormalTok{reo\_start\_mask }\OperatorTok{=}\NormalTok{ (}\OperatorTok{\textasciitilde{}}\NormalTok{is\_reo\_padded[:}\OperatorTok{{-}}\DecValTok{1}\NormalTok{]) }\OperatorTok{\&}\NormalTok{ is\_reo\_padded[}\DecValTok{1}\NormalTok{:]}
\NormalTok{is\_reorientation\_start[reo\_start\_mask] }\OperatorTok{=} \VariableTok{True}

\CommentTok{\# Count only starts:}
\NormalTok{n\_reorientations }\OperatorTok{=}\NormalTok{ is\_reorientation\_start.}\BuiltInTok{sum}\NormalTok{()  }\CommentTok{\# Correct}
\CommentTok{\# NOT: n\_reorientations = is\_reorientation.sum()  \# Wrong {-} counts duration!}
\end{Highlighting}
\end{Shaded}

\hypertarget{turn-definition}{%
\subsection{Turn Definition}\label{turn-definition}}

\textbf{CRITICAL}: Turns are reorientations that contain pauses. This is
a behavioral requirement.

\begin{Shaded}
\begin{Highlighting}[]
\CommentTok{\# Correct approach:}
\ControlFlowTok{for}\NormalTok{ reo\_start, reo\_end }\KeywordTok{in}\NormalTok{ reorientations:}
\NormalTok{    reo\_pause\_frames }\OperatorTok{=}\NormalTok{ is\_pause[reo\_start:reo\_end}\OperatorTok{+}\DecValTok{1}\NormalTok{]}
    \ControlFlowTok{if}\NormalTok{ np.}\BuiltInTok{any}\NormalTok{(reo\_pause\_frames):}
\NormalTok{        turns.append((reo\_start, reo\_end))  }\CommentTok{\# This is a turn}
\end{Highlighting}
\end{Shaded}

\hypertarget{cycle-extraction}{%
\subsection{Cycle Extraction}\label{cycle-extraction}}

Cycles are defined relative to stimulus onset (t=0), with: - Baseline
period: 10 seconds before onset - Pulse period: from onset to LED OFF -
Analysis window: baseline + pulse

\hypertarget{data-structures}{%
\section{Data Structures}\label{data-structures}}

\hypertarget{segmentation-dictionary}{%
\subsection{Segmentation Dictionary}\label{segmentation-dictionary}}

\begin{Shaded}
\begin{Highlighting}[]
\NormalTok{segmentation }\OperatorTok{=}\NormalTok{ \{}
    \StringTok{\textquotesingle{}runs\textquotesingle{}}\NormalTok{: [(start\_idx, end\_idx), ...],}
    \StringTok{\textquotesingle{}head\_swings\textquotesingle{}}\NormalTok{: [(start\_idx, end\_idx), ...],}
    \StringTok{\textquotesingle{}reorientations\textquotesingle{}}\NormalTok{: [(start\_idx, end\_idx), ...],}
    \StringTok{\textquotesingle{}turns\textquotesingle{}}\NormalTok{: [(start\_idx, end\_idx), ...],  }\CommentTok{\# NEW: reorientations with pauses}
    \StringTok{\textquotesingle{}is\_run\textquotesingle{}}\NormalTok{: np.array([}\BuiltInTok{bool}\NormalTok{, ...]),}
    \StringTok{\textquotesingle{}is\_head\_swing\textquotesingle{}}\NormalTok{: np.array([}\BuiltInTok{bool}\NormalTok{, ...]),}
    \StringTok{\textquotesingle{}is\_reorientation\textquotesingle{}}\NormalTok{: np.array([}\BuiltInTok{bool}\NormalTok{, ...]),  }\CommentTok{\# Start events only}
    \StringTok{\textquotesingle{}is\_turn\textquotesingle{}}\NormalTok{: np.array([}\BuiltInTok{bool}\NormalTok{, ...]),  }\CommentTok{\# NEW: turns (reorientations with pauses)}
    \StringTok{\textquotesingle{}n\_runs\textquotesingle{}}\NormalTok{: }\BuiltInTok{int}\NormalTok{,}
    \StringTok{\textquotesingle{}n\_head\_swings\textquotesingle{}}\NormalTok{: }\BuiltInTok{int}\NormalTok{,}
    \StringTok{\textquotesingle{}n\_reorientations\textquotesingle{}}\NormalTok{: }\BuiltInTok{int}\NormalTok{,}
    \StringTok{\textquotesingle{}n\_turns\textquotesingle{}}\NormalTok{: }\BuiltInTok{int}  \CommentTok{\# NEW}
\NormalTok{\}}
\end{Highlighting}
\end{Shaded}

\hypertarget{cycle-dictionary}{%
\subsection{Cycle Dictionary}\label{cycle-dictionary}}

\begin{Shaded}
\begin{Highlighting}[]
\NormalTok{cycle }\OperatorTok{=}\NormalTok{ \{}
    \StringTok{\textquotesingle{}cycle\_num\textquotesingle{}}\NormalTok{: }\BuiltInTok{int}\NormalTok{,}
    \StringTok{\textquotesingle{}onset\_frame\textquotesingle{}}\NormalTok{: }\BuiltInTok{int}\NormalTok{,  }\CommentTok{\# Stimulus onset frame (t=0 reference)}
    \StringTok{\textquotesingle{}onset\_time\textquotesingle{}}\NormalTok{: }\BuiltInTok{float}\NormalTok{,  }\CommentTok{\# Stimulus onset time (t=0)}
    \StringTok{\textquotesingle{}led\_on\_frame\textquotesingle{}}\NormalTok{: }\BuiltInTok{int}\NormalTok{,}
    \StringTok{\textquotesingle{}led\_on\_time\textquotesingle{}}\NormalTok{: }\BuiltInTok{float}\NormalTok{,}
    \StringTok{\textquotesingle{}led\_off\_frame\textquotesingle{}}\NormalTok{: }\BuiltInTok{int}\NormalTok{,}
    \StringTok{\textquotesingle{}led\_off\_time\textquotesingle{}}\NormalTok{: }\BuiltInTok{float}\NormalTok{,}
    \StringTok{\textquotesingle{}pulse\_duration\textquotesingle{}}\NormalTok{: }\BuiltInTok{float}\NormalTok{,}
    \StringTok{\textquotesingle{}eti\textquotesingle{}}\NormalTok{: }\BuiltInTok{float}\NormalTok{,  }\CommentTok{\# Elapsed time index (OFF to next onset)}
    \StringTok{\textquotesingle{}baseline\_start\_time\textquotesingle{}}\NormalTok{: }\BuiltInTok{float}\NormalTok{,  }\CommentTok{\# 10s before onset}
    \StringTok{\textquotesingle{}cycle\_start\_time\textquotesingle{}}\NormalTok{: }\BuiltInTok{float}\NormalTok{,  }\CommentTok{\# Start of analysis window}
    \StringTok{\textquotesingle{}cycle\_end\_time\textquotesingle{}}\NormalTok{: }\BuiltInTok{float}  \CommentTok{\# End of analysis window (LED OFF)}
\NormalTok{\}}
\end{Highlighting}
\end{Shaded}

\hypertarget{references}{%
\section{References}\label{references}}

\begin{itemize}
\tightlist
\item
  MAGAT Segmentation: Based on \texttt{@MaggotTrack/segmentTrack.m} from
  magniphyq repository
\item
  Klein Run Table: Based on Mason Klein's 18-column run table
  methodology
\item
  Turn Rate Calculation: Matches MATLAB \texttt{ExperimentAggregator.m}
  methodology
\end{itemize}



\end{document}
